\chapter{Debugging on the Chip}
\label{ch:debug}

When something misbehaves on the silicon there are three complementary tools: an
interactive debugger over JTAG, the post-mortem backtrace, and the runtime's own
fault diagnostics. Each works against the bare application image directly, using
only the pinned cross-toolchain.

\section{Interactive debugging: OpenOCD + GDB}\index{debugging}\index{OpenOCD}\index{GDB}\index{JTAG}\index{USB-JTAG}\index{breakpoint}\index{backtrace}

The ESP32-S3's built-in USB-Serial-JTAG port is both the console \emph{and} a JTAG
debug interface, so a single USB cable is enough. The project pins its own copies
of the debug tools; one command fetches them:

\begin{lstlisting}[style=sh]
./x get-debug-tools      # downloads OpenOCD + xtensa-esp32s3-elf-gdb into tools/
./x debug blink          # build, flash, start OpenOCD, attach GDB, rest at Main
\end{lstlisting}

One subtlety is worth knowing: the S3 needs the \emph{S3-specific}
\texttt{xtensa-esp32s3-elf-gdb}; the Alire crate's esp32 (LX6) GDB silently fails
on it. \texttt{./x debug} selects the right one. From there it is ordinary GDB on
\texttt{app.elf}: breakpoints, \texttt{step}, \texttt{print},
\texttt{backtrace}. Two options matter on this target:

\begin{itemize}
  \item \texttt{--smp} exposes \emph{both} LX7 cores as GDB threads
        (\texttt{info threads}), so you can see what each core is doing,
        essential for a cross-core hang.
  \item \texttt{--attach} halts the board \emph{in place} without resetting it, so
        you can examine a live hang or crash exactly where it stopped (a
        post-mortem).
\end{itemize}

The command line and both editors (next section) drive this same OpenOCD + GDB
path; when finished, \texttt{./x kill-openocd} releases the USB-JTAG port.

\section{Editor integration: VS Code and Vim}\index{editor integration}\index{Ada Language Server}

The same build, flash and debug actions are wired into both VS Code and Vim, so
you need never leave the editor. Both drive the top-level \texttt{./x} dispatcher
and both use the project's pinned OpenOCD + s3 GDB for on-chip debugging.

Two prerequisites are shared by either editor. On-chip debugging first needs the
tools fetched once:

\begin{lstlisting}[style=sh]
./x get-debug-tools     # pinned, SHA-256-verified OpenOCD + xtensa-esp32s3-elf-gdb
\end{lstlisting}

And \emph{code intelligence} (completion, go-to-definition, diagnostics, hover)
comes from the \textbf{Ada Language Server}, which ships inside the AdaCore
toolchain (nothing to install separately). It only needs to be pointed at the
example's \texttt{.gpr} project file, shown for each editor below.

\subsection{VS Code}\index{VS Code}

There are two complementary pieces. The first is the committed workspace
configuration in \texttt{.vscode/} (\texttt{tasks.json}, \texttt{launch.json},
\texttt{settings.json}), which wires the editor to \texttt{./x} and needs two
extensions:

\begin{description}[style=nextline]
\item[Ada \& SPARK (\texttt{AdaCore.ada}).] The Ada Language Server: required,
  for all code intelligence.
\item[Native Debug (\texttt{webfreak.debug}).] A GDB frontend, needed only for
  \texttt{F5} debugging. \texttt{cppdbg} cannot walk the Xtensa register-window
  stack (it connects but shows no call stack or source); Native Debug simply
  relays GDB's own frames, which unwind correctly.
\end{description}

Install both from a terminal (or from the Extensions view, which also offers them
as workspace recommendations):

\begin{lstlisting}[style=sh]
code --install-extension AdaCore.ada
code --install-extension webfreak.debug
\end{lstlisting}

Then point the language server at the example you are working on by editing
\texttt{ada.projectFile} in \texttt{.vscode/settings.json} and reloading the
window (\texttt{Ctrl-Shift-P} $\rightarrow$ \emph{Developer: Reload Window}):

\begin{lstlisting}[style=sh]
"ada.projectFile": "examples/esp32s3_gpio0_blink/gpio_blink.gpr"
\end{lstlisting}

Build, flash and monitor are tasks: \texttt{Ctrl-Shift-B} runs the default
\texttt{Ada: build}, and \emph{Tasks: Run Task} offers flash, monitor, run and
clean. Each prompts for the example, serial port and runtime profile and routes
to \texttt{./x}; build errors land in the \textbf{Problems} panel via a GNAT
problem matcher.

To \textbf{debug} (\texttt{F5}, ``Ada: debug (OpenOCD + GDB)''): build \emph{and
flash} the example first, so the chip runs the same image whose \texttt{app.elf}
you debug. The launch only \emph{attaches}, so a stale or different-profile
image leaves the breakpoint addresses mismatched. It starts OpenOCD and halts at
\texttt{app\_main} (the shared boot-glue entry; these GNAT images have no C
\texttt{main}); set breakpoints in the Ada source and continue/step.

The second piece is an optional \textbf{status-bar extension}
(\texttt{ide/vscode-ada-esp32}) that adds clickable build / flash / monitor /
debug icons plus persistent port and profile selectors. A prebuilt
\texttt{.vsix} is committed to the repository, so installing needs only the
\texttt{code} CLI on \texttt{PATH}:

\begin{lstlisting}[style=sh]
./x install-ide         # installs the committed vscode-ada-esp32.vsix
\end{lstlisting}

Reload the window afterwards; re-run it to update after a \texttt{git pull} (a
repo-hosted \texttt{.vsix} does not auto-update). Its debug icon uses the same
OpenOCD + GDB launch described above.

\subsection{Vim and Neovim}\index{Vim}\index{termdebug}

The \texttt{ide/vim-ada-esp32} plugin gives the same actions as Vim commands. It
is pure Vimscript (nothing to compile) and needs \textbf{Vim 8.1+ with
\texttt{+terminal}} (which bundles the \texttt{termdebug} GDB frontend); Neovim
works too. Install it with one command, which symlinks the plugin into Vim's and
Neovim's native package directories; because it is a symlink \emph{into} the repo,
a \texttt{git pull} updates it in place:

\begin{lstlisting}[style=sh]
./x install-vim
\end{lstlisting}

Or wire it up by hand; it has the standard \texttt{plugin/} layout:

\begin{lstlisting}[style=sh]
" vim-plug
Plug '/path/to/ada_esp32s3/ide/vim-ada-esp32'
" or, with no plugin manager:
set runtimepath^=/path/to/ada_esp32s3/ide/vim-ada-esp32
\end{lstlisting}

The commands auto-detect context by searching upward for \texttt{./x}, so they
work from anywhere inside the repo:

\begin{description}[style=nextline]
\item[{\texttt{:AdaEsp32Build [example]}}] build via \texttt{./x build}.
\item[{\texttt{:AdaEsp32Flash [example]}}] flash to the selected port.
\item[{\texttt{:AdaEsp32BuildFlash [example]}}] build then flash.
\item[{\texttt{:AdaEsp32Run [example]}}] build $+$ flash $+$ monitor.
\item[{\texttt{:AdaEsp32Monitor}}] serial console at 115200.
\item[{\texttt{:AdaEsp32Debug [example]}}] OpenOCD $+$ \texttt{termdebug} (below).
\item[{\texttt{:AdaEsp32Port [port]} / \texttt{:AdaEsp32Profile [prof]}}] set the
  serial port / runtime profile; both persist across later commands.
\end{description}

The example argument accepts the short name (\texttt{gpio0\_blink}) or the full
directory (\texttt{esp32s3\_gpio0\_blink}), with tab-completion; the chosen
example, port and profile persist (in \texttt{g:ada\_esp32\_example} and friends)
so later commands reuse them without re-prompting.

\textbf{Debugging} uses Vim's built-in \texttt{termdebug}, the Vim analog of
the Native Debug setup, relaying GDB's own frames so it handles Xtensa correctly.
After a build:

\begin{lstlisting}[style=sh]
:AdaEsp32Build gpio0_blink
:AdaEsp32Debug gpio0_blink
\end{lstlisting}

This starts OpenOCD in the background, waits for its GDB server, then opens
\texttt{termdebug} reset and halted with a breakpoint armed at \texttt{app\_main}.
From there \texttt{:Continue} runs to it, \texttt{:Break src/foo.adb:53} (or
\texttt{:Break} on the cursor line) sets an Ada breakpoint, and
\texttt{:Step}/\texttt{:Over}/\texttt{:Finish} are the step controls; OpenOCD is
stopped automatically when you quit the session. One inherent caveat: a
\texttt{:Over} across a \emph{blocking} statement (a \texttt{delay}, an entry
call) can step into the hand-written Xtensa assembly (the interrupt vectors or
context switch), which carries no debug bounds. GDB then reports ``cannot find
bounds of current function.'' Set a breakpoint past the statement and
\texttt{:Continue} instead.

Code intelligence again comes from the Ada Language Server through your LSP
client, pointed at the example's \texttt{.gpr}. For example, with
\texttt{coc.nvim} in \texttt{coc-settings.json}:

\begin{lstlisting}[style=sh]
{ "ada.projectFile": "examples/esp32s3_gpio0_blink/gpio_blink.gpr" }
\end{lstlisting}

or register \texttt{ada\_language\_server} for the \texttt{ada} filetype with
native LSP / vim-lsp and set the same \texttt{ada.projectFile} setting.

\section{Post-mortem: decoding a Guru Meditation}\index{Guru Meditation}\index{panic}\index{addr2line}\index{post-mortem}

An unhandled fault prints a register dump and a backtrace of program-counter
values over the console (a ``Guru Meditation''). You do not need a live debugger
to read it: turn the addresses back into source with \texttt{addr2line} against
the ELF:

\begin{lstlisting}[style=sh]
xtensa-esp32s3-elf-addr2line -fe build/app.elf 0x4200abcd 0x4200ef01
\end{lstlisting}

In the dump, \texttt{EXCVADDR} is the faulting data address and \texttt{EXCCAUSE}
the reason: a \texttt{LoadStoreError} (a bad data access), or an
\texttt{InstructionFetchError} (a wild jump, often a corrupted stack or a stray
pointer). Flash, reproduce, decode: that loop localises most faults without ever
attaching a debugger.

\section{The runtime's own diagnostics}\index{stack overflow}\index{last-chance handler}\index{safe state}

The runtime helps before you reach for a tool at all:

\begin{itemize}
  \item \textbf{Precise stack-overflow detection.} Each task arms a hardware
        watchpoint near its stack limit, so an overflow faults \emph{precisely}
        (a debug exception) instead of silently corrupting a neighbour, and on
        the \texttt{full} profile the runtime can turn it into a catchable
        \texttt{Storage\_Error}.
  \item \textbf{A fault safe-state hook.} The \texttt{esp32s3\_exceptions}
        example wraps the panic handler to classify a fault (stack overflow
        versus other) against the running task's stack window, latch a safe
        state, and print an explicit \texttt{*** STACK OVERFLOW ***} verdict
        before the board resets.
  \item \textbf{The last-chance handler.} Any exception that escapes every handler
        ends up in the runtime's last-chance handler; on the leanest profile,
        where exceptions cannot propagate across calls, that is your fault report:
        a message and a controlled reset.
\end{itemize}

\textbf{A note on console output while debugging.} \texttt{Ada.Text\_IO} reaches
\texttt{/dev/ttyACM0}; for terse status from interrupt or early-boot context the
examples call the ROM \texttt{esp\_rom\_printf}, which is always available. When
two cores (or an interrupt handler and a task) might print at once, funnel
output through one task: interleaved single-character writes have themselves been
mistaken for bugs more than once.

\section{Breadcrumbs at the panic park}\index{panic}\index{breadcrumbs}\index{double exception}\index{PRID}

When a fault escalates past every handler \textemdash{} an unhandled
exception, a double exception, a debug exception with no debugger attached
\textemdash{} the vendored vectors funnel into \texttt{\_xt\_panic}, which
used to be a bare \texttt{j .}: the board parked with the evidence locked in
special registers that the next reset destroys.  \texttt{start.S} now dumps
the machine state to a fixed block (\texttt{\_\_gnat\_panic\_dump}) before
parking: EXCCAUSE and EXCVADDR, EPC1..EPC7 with their EPS twins, DEPC, PS,
\texttt{a0}/\texttt{a1}, WINDOWBASE/WINDOWSTART, PRID, and a marker for which
entry point was taken.  The first writer wins \textemdash{} a magic word gates
the block, so when both cores pile into the park the original fault survives
\textemdash{} and the block lives in \texttt{.data}, re-initialised at every
boot, so a stale dump cannot masquerade as a fresh one.  One JTAG
\texttt{mdw} reads the whole story out of a wedged board.

Two decoding traps, both paid for in hours.  The double-exception vector
\emph{overwrites} EXCCAUSE with a panic-reason code before calling
\texttt{\_xt\_panic}, so ``cause 2'' at a double-exception park means
\emph{double exception}, not \emph{instruction fetch error}.  And PRID does
not name the core by its whole value: bit~13 alone does.  On the ESP32-S3,
\texttt{16\#CDCD\#} is core~0 and \texttt{16\#ABAB\#} is core~1
\textemdash{} misread that and every per-core deduction that follows is
backwards.

\section{Hardware watchpoints without a debugger}\index{DBREAK}\index{watchpoint}\index{debug exception}

OpenOCD watchpoints have two failure modes on this chip: they cannot be armed
early enough to catch a corruption that happens inside the first second of
boot, and an armed watchpoint over an address the ROM itself uses turns reset
into a watchdog loop.  The escape is to let the \emph{firmware} arm the
watchpoint at exactly the right moment \textemdash{} three instructions:

\begin{lstlisting}[style=sh]
wsr.dbreaka0  <address>       # the word to watch
wsr.dbreakc0  16#8000_003C#   # store-break, 4-byte mask
isync                          # (16#4000_003C# = load-break)
\end{lstlisting}

If a debugger is attached, the hit halts the core at the offending store with
the writer's PC on screen.  If \emph{no} debugger is attached the hit becomes
a debug exception that funnels into \texttt{\_xt\_panic} \textemdash{} and
the breadcrumb block's EPC6 slot (the debug level is~6) records the writer's
PC anyway.  A store-watch armed after setup, plus one \texttt{mdw} after the
park, identified in one boot a rogue writer that three days of halt-and-inspect
archaeology had missed.

The technique has a sting worth knowing even if you never use it: DBREAK state
\emph{survives soft resets}.  The full runtime's stack-overflow guard arms a
DBREAK in normal operation, so reflashing from a full-profile image to any
other image used to inherit a live watchpoint over what is now unrelated data
\textemdash{} and the new image died at boot with a debuggerless debug
exception.  \texttt{start.S} therefore zeroes \texttt{DBREAKC0/1} and
\texttt{IBREAKENABLE} on both cores before anything else runs; if you port the
boot path elsewhere, keep that.

\section{Decoding the compiler's bare-metal complaints}
\index{No\_Implicit\_Dynamic\_Code}\index{accessibility check}\index{restrictions}

The bare-metal profiles forbid things a hosted program takes for granted, and the
compiler enforces the bans. The messages are precise but terse, and the first time
one stops a build it can be baffling. Here are the ones you are most likely to meet
coming from C, what each really means, and the fix.

\paragraph{\texttt{subprogram must not be deeper than access type}.}
You took \texttt{'Access} of a \emph{nested} subprogram (one declared inside
another subprogram) to use as a callback. On a hosted target the compiler would
build a \emph{trampoline} (a scrap of code on the stack) to carry the enclosing
frame; the \texttt{No\_Implicit\_Dynamic\_Code} restriction forbids writing code to
the stack at all, and on this chip such a trampoline faults outright. \emph{Fix:}
move the callback to library level (the top of a package) and pass any state it
needs through a context parameter rather than capturing it: closure-free. This is
why every callback in this book is a library-level, \texttt{Ctx}-carrying
subprogram (\S\ref{ch:ada-primer} introduces subprograms; the HAL callbacks follow
this rule throughout).

\paragraph{\texttt{accessibility check failed} (a \texttt{Program\_Error} at run
time), or \texttt{non-local pointer cannot point to local object} (at compile
time).} You stored an \texttt{'Access} of an object into something that outlives it
into a library-level variable, a registry, a longer-lived record. Once the object's
scope ends the stored pointer would dangle, so Ada rejects it: at compile time when
it can prove the lifetimes, at run time otherwise. \emph{Fix:} give the object the
lifetime the pointer assumes, which usually means declaring it at library level.
(This is exactly the bug behind the FTP server that stored a \texttt{Main}-local
filesystem mount in a library-level handler, cured by moving the mount into a
library-level package.)

\paragraph{\texttt{violation of restriction No\_X}.}
You used a feature the chosen runtime profile bans: a heap allocator, a second
task entry, an exception handler under a no-propagation profile, and so on. The
profile is a contract, and the compiler holds you to it. \emph{Fix:} either avoid
the feature, or build under a richer profile (light-tasking $\rightarrow$ embedded
$\rightarrow$ full) that permits it; Chapter~\ref{ch:profiles} maps features to
profiles.

\paragraph{\texttt{Constraint\_Error} (a range or validity check).}
A value left the range its (sub)type allows: an index past the end of an array,
an out-of-range assignment, a bad enumeration read from external data. Unlike the
above, this is almost always a genuine bug in your data or arithmetic rather than a
profile rule. \emph{Fix:} find the offending value; the type told you its range, so
the check is doing its job. Build with checks on (\texttt{-gnata}) while developing;
a release build may have them off, turning the same mistake into silent
corruption.

The throughline is that each message names a guarantee the bare-metal target makes
(no stack-written code, no dangling pointers, no banned features, no out-of-range
values) and points at the one place you broke it. They read as obstacles at first
and as a safety net soon after.
